\documentclass{jsarticle}
\usepackage{amsmath}
\usepackage{amssymb}
\begin{document}
\title{Ordinary differential manifold deconstinuace with \\
global manifold of estimate meanings.}
\author{Masaaki Yamaguchi}
\date{}
\maketitle

    $$ \int f(x)dx = \int (4x^5 + x^4 + 3x^2 + x) dx$$
    この積分記号の中で微分をすると、
    $$ \int f(x)^{'} dx = \int ({20x^4 + 4x^3 + 6x +1})dx $$
    $$ = 4x^5 + x^4 + 3x^2 + x + \text{C(Cは積分記号)} $$
    歴然と、外の積分記号が中の微分で消えている。
    常微分と常積分では、互いに記号が相殺される。
    ファインマン博士が、積分記号の中で微分してみたらは、
    部分積分多様体のことを言っているのは、歴然としている。
    $$ \int f(x)^{'}dx = [xF(x)] - \int x^{'}f(x)dx $$
    $$ = [xF(x)] - F(x) + \text{C(Cは積分定数)} $$
    これは、大域的積分多様体の質量差エネルギーの式になっている。
    一般相対性理論と特殊相対性理論での$ E = mc^2 - {1 \over 2}mv^2 $
    となっている。
    $$ = \int \Gamma(x)dx $$
    $$ = x\Gamma(x) $$

    $$ \int \Gamma(\gamma)^{'}dx_m = 2e^{-x \log x} $$
    $$ T =  \left(\int \Gamma(\gamma)^{'}dx_m\right)^{'}   $$
    ガンマ関数の大域的積分多様体を常微分多様体で解くと、
    フェルマーの定理になる。
     $$ = -2 e^{-x \log x} $$
    ガンマ関数の大域的積分多様体に大域的微分作用素を施すと、
    $$ {(T^{\nabla})}^{\oplus L} = e^{-f} + e^{f} \ge e $$
    これは、相加相乗平均の方程式になる。
    $$ f(x) + g(x) \ge {2\sqrt{f(x)\cdot g(x)}} $$
    要するに、常微分方程式になるのを、大域的作用素では、
    エターナルな空間になる。曲平面になる。
    微分では、平面から点への極限が、大域的では、点から平面になり、
    だが、ガンマ関数の大域的２重微分多様体は、特異点を求めることになり、
    大域的積分は、点から平面、立方体を求めることになる。
    説明が矛盾しているように聞こえるが、M理論の絡み合いの
    種数を、平均場理論で求めていることということである。
    なぜかは、
    $$ x^x = e^{x \log x} $$
    を使うからであり、そのようになっている理論が、M理論である。
    大域的微分作用素で、点と世界線が同型と求まっている。
    閉３次元多様体と特異点が同型となり、種数が違うのが、どうしてかが、
    この方程式で求まってもいる。

    フェルマーの定理は、２通りの解がある。
    １つ目は、$ -2e^{x \log x} $であり、２つ目は、$ -2 e^{- x \log x} $
    である。
    不確定性原理の位置と運動量での、宇宙と異次元の関係が、
    フェルマーの定理であり、ゼータ関数の言葉の意味そのままである。
    その解の構成が、Jones多項式である。
    大域的部分積分多様体と、一般の部分積分多様体とは、
    積分の中で微分すると部分積分多様体になり、
    両方共に、常積分多様体とは
    全然違う解き方になる。
    $$ (\log x)^{'} = {1 \over x}, x^n+ y^n = z^n $$
    $$ x^n = -y^n + c, nx^{n-1} = -ny^{n-1}y^{'} $$
    $$ y^{'} ={{nx^{n-1}} \over {ny^{n-1}}} $$
    $$ ={x^{n-1} \over y^{n-1}} = - {{y \over x}\cdot{x \over y}^n} $$
    $$ -{{\cos x} \over (\cos x)^{'}}(\sin x)^{'} = z_n $$
    $$ z^n = - 2 e^{x \log x} $$
    $$ \int C dx_m = {d \over df}F(x) $$
    これは、上の常識を超える計算である。
    積分の外で微分しても、中の積分が消えない。
    積分の中で微分しても、積分が消えない。
    大域的微分と大域的積分多様体は、革命的な微分と多様体積分である。
    $$  \int \Gamma(\gamma)^{'}dx_m = \int \Gamma dx_m \cdot {d \over 
    d \gamma}\Gamma \le \int \Gamma dx_m + {d \over d \gamma}\Gamma $$
    $$ = \Gamma^{\Gamma} + {\Gamma^{\gamma}}^{'} $$
    $$  \int C dx_m = {d \over df}F(x) = \int(\int {1 \over x^s}dx - \log x)
    \mathrm{dvol} $$
    $$ = {d \over df}\int \int{1 \over {(y \log y)^{1 \over 2}}}dy_m - 
    \int e^x x^{1 - t}dx_m $$
    $$ = {d \over df}\int\int{1 \over {(x \log x)^2}}dx_m +
    {d \over df}\int\int{1 \over {(y \log y)^{1 \over 2}}}dy_m$$
    大域的多様体は、革命的な微分と積分の計算方法を、私は、
    アカシックレコードの子に伝授された。
   $$ i = (1,0)\cdot (1,0), e^{i\theta} = \cos \theta + i \sin \theta $$
   $$ e^{-i \theta} = \cos \theta - i\sin \theta $$
   $$ \sin \theta = {e^{i \theta} - e^{-i \theta} \over 2i} $$
   $$ \sin i \theta = {e^{-\theta} - e^{\theta} \over 2} $$
   $$ \pi (\chi,x) = \cos \theta + i \sin \theta $$

   In this equations, two dimension redestructed into three dimension,
   this destroy of reconstructed way is append with fifth dimension.
   This deconstructed way of redestructe is arround of universe attached
   with three dimension, this over cover call into fifth dimension.

   $$ R(-\alpha) = \begin{pmatrix} \cos \alpha & \sin \alpha \\
   - \sin \alpha & \cos \alpha \end{pmatrix} $$
   $$ R(\alpha) = \begin{pmatrix} \cos \alpha & - \sin \alpha \\
   \sin \alpha & \cos \alpha \end{pmatrix} $$
   $$ R(\alpha)M R(-\alpha) = 
   \begin{pmatrix} \cos \alpha & - \sin \alpha \\
   \sin \alpha & \cos \alpha \end{pmatrix} \begin{pmatrix} 1 & 0 \\
   0 & - 1 \end{pmatrix} \begin{pmatrix} \cos \alpha & \sin \alpha \\
   - \sin \alpha & \cos \alpha \end{pmatrix} $$
   $$ = \begin{pmatrix} \cos^2 \alpha - \sin^2 \alpha & 2 \sin \alpha \cos 
	   \alpha \\
	   2 \sin \alpha \cos \alpha & - \cos^2 \alpha + \sin^2 \alpha
   \end{pmatrix} $$
   $$ = \begin{pmatrix} \cos 2 \alpha & \sin 2 \alpha \\
   \sin 2 \alpha & - \cos 2 \alpha \end{pmatrix} $$

   $$ \begin{pmatrix} \cos \alpha & - 1 \\
	   1 & - \sin \alpha \end{pmatrix} \begin{pmatrix} \cos \alpha \\ 
		   \sin \alpha
   \end{pmatrix} = \begin{pmatrix} 1 & 0 \\
   0 & - 1 \end{pmatrix} $$
   $$ \begin{pmatrix} \cos \alpha & \sin \alpha \\
   \sin \alpha & - \cos \alpha \end{pmatrix} \begin{pmatrix}
   x \\ y \end{pmatrix} = \begin{pmatrix} 1 & 0 \\
   0 & - 1 \end{pmatrix} $$
   $$ \lim_{\theta \to 0} \begin{pmatrix} \sin \theta \\ \cos \theta 
   \end{pmatrix} \begin{pmatrix} \theta & 1 \\
   1 & \theta \end{pmatrix} \begin{pmatrix} \cos \theta \\ \sin \theta
	   \end{pmatrix} = \begin{pmatrix} 1 & 0 \\ 0 & - 1 \end{pmatrix} $$
		   $$ f^{-1}xf(x) = 1 $$
		   $$ (\log x)^{'} = {1 \over x}, x^n + y^n = z^n  $$
		   $$ x^n = - y^n + c, nx^{n-1} = -n y^{n-1}y^{'} $$
		   $$ y^{'} = {{n x^{n-1}} \over {n y^{n-1}}} $$
		   $$ = {x^{n-1} \over y^{n-1}}= - {y \over x}\cdot
		   ({x \over y})^n $$
		   $$ - {\cos x \over (\cos x)^{'}}(\sin x)^{'} = z_n $$
		   $$ z^n = - 2 e^{x \log x} $$
		   $$ \lim_{x \to \infty} f(x) = a, \lim_{y \to \infty}
		   f(y) = b, \lim_{x,y \to \infty}\{f(x) + f(y)\} = a + b  $$
		   $$ \lim_{z \to 0}f(z) = c, \delta \int z^n = {d \over dV}
		   x^3 $$
		   $$ \lim_{x \to \infty} = c - \lim_{y \to \infty}f(y)
		   \lim_{x \to \infty} f(x) + \lim_{y \to \infty}f(y) = 
		   \lim_{z \to \infty}f(z) $$
		   $$ z^n = \cos n \theta + i \sin n \theta $$
		   $$ = - 2 e^{x \log x} $$

  These equation is gravity and antigravity equation. 

	   $$ {d \over d\sigma}\begin{bmatrix}{(\sigma_1 + \sigma_2) 
	   \over 2}\end{bmatrix}  $$
	   $$ = \sigma(\upharpoonleft \downharpoonleft) + \sigma(
	   \upuparrows)+ \sigma(\downdownarrows)+ 
	   \sigma(\leftleftarrows)
	   + \sigma(\rightrightarrows) $$
	   $$ {d \over df}{\int \int {1 \over (x \log x)^2}}dx_m
	   = \sigma(\upharpoonright) $$
	   $$ {d \over df}{\int \int {1 \over (y \log y)^{1 \over2}}}dy_m
	   = \sigma(\downharpoonright) $$
	   $$ \sigma(\leftarrowtail) + \sigma(\rightarrowtail)
	   = \int e^{-f}[-\Delta v + R_{ij}v_{ij} + \nabla_{i}\nabla_{j}v
	   + v\nabla_{i}\nabla_{j} + 2 <f,h> + (R + \nabla f)(v - 
	   {h \over 2})]
	   $$
	   $$ \sigma(\downharpoonleft) = \sigma(\upharpoonright 
	   \downharpoonright + \upuparrows 
	   + \rightrightarrows)$$
	   $$ \sigma(\upharpoonleft) =  \sigma(\upharpoonright 
	   \downharpoonright + \downdownarrows 
	   + \leftleftarrows) $$
         
	   $$ \text{weak electric theorem} = \sigma(\upharpoonleft) $$
	   $$ \text{strong electric theorem} = \sigma(\downharpoonleft) $$

	   These equation is represented with topology of string model,
	   and weak electric theorem is constructed with Maxwell theorem 
	   and weak boson, gravity that estimate with this three power united.
	   Moreover, strong boson and Maxwell theorem, antigravity that
	   also estimate with this three power united.
	   This two united power is integrated from gravity and antigravity.
	   Then this united power is zeta function.

	    
	   $$ y =  1 + \tan^2 x = {1 \over {\cos x}} $$
	   $$ y^{'} = {{\cos^{'}x} \over {\cos^2 x}} $$
           $$ -{{\cos x} \over (\cos x)^{'}}(\sin x)^{'} = z_n $$
	   $$ {1 \over y^{'}} = z^2_n$$
           $$ z^n = - 2 e^{x \log x} $$

	   ガンマ関数の大域的積分多様体は、結局は、フェルマーの定理
	   だった。

         代数的計算手法のために$ \oplus L $を使っている。
そのために、冪乗計算と商代数の計算が、乗算で楽に見えるようになっている。
            微分幾何の量子化は、代数幾何の量子化の計算になっている。
  加減乗除が初等幾何であり、大域的微分と積分が、現代数学の代数計算の
  簡易での楽になる計算になっている。
  初等代数の計算は、
  $$ \bigoplus({i \hbar^{\nabla}})^{\oplus L} $$
   $$m \bigoplus({i \hbar^{\nabla}})^{\oplus L} +  
   n\bigoplus({i \hbar^{\nabla}})^{\oplus L} = 
   (m+n) \bigoplus({i \hbar^{\nabla}})^{\oplus L}   $$
   $$ m\bigoplus({i \hbar^{\nabla}})^{\oplus L} -  
   n\bigoplus({i \hbar^{\nabla}})^{\oplus L} = 
   (m-n) \bigoplus({i \hbar^{\nabla}})^{\oplus L}   $$
   $$ \bigoplus({i \hbar^{\nabla}})^{{\oplus L}^m} \times 
      \bigoplus({i \hbar^{\nabla}})^{{\oplus L}^n} =
   \bigoplus({i \hbar^{\nabla}})^{\oplus L^{m+n}} $$
  $$ {{\bigoplus({i \hbar^{\nabla}})^{\oplus L^m}}  \over  
  {\bigoplus({i \hbar^{\nabla}})^{\oplus L^n}}} =  
  \bigoplus({i \hbar^{\nabla}})^{\oplus L^{m\over n}} $$
  大域的計算での微分と積分は、
  $$ \left(\bigoplus({i \hbar^{\nabla}})^{\oplus L}\right)^{df}
  =  {\left(\bigoplus({i \hbar^{\nabla}})^{\oplus L}\right)^{
	  \bigoplus({i \hbar^{\nabla}})^{\oplus L}}}^{'}  $$
  $$ = \bigoplus({i \hbar^{\nabla}})^{\oplus L^{'}} $$
   $$ \int \bigoplus({i \hbar^{\nabla}})^{\oplus L} dx_m $$
   $$ \bigoplus({i \hbar^{\nabla}})^{\oplus L} = n $$
   大域的積分の代数幾何の量子化での計算は、ガンマ関数になっている。
   $$ {n^{L+1} \over {n+1}} = \int e^x x^{1 - t} dx $$
   代数幾何の量子化の因数分解は、加減乗除の式のm,nの
   組み合わせ多様体でのガンマ関数同士の計算になる。
   $$ \left(\bigoplus({i \hbar^{\nabla}})^{\oplus L}+m\right) 
   \left(\bigoplus({i \hbar^{\nabla}})^{\oplus L}+n\right) $$
  $$ = {n^{L+1} \over {n+1}} = t \int (1 - x)^n x^{m - 1} dx $$
  $$ \int x^{m-1}(1 - x)^{n-1} dx $$
  $$ \nabla ({i \hbar^{\nabla}})^{\oplus L}, \bigoplus ({i \hbar^{\nabla}})^{
	   \oplus L}, \square ({i \hbar^{\nabla}})^{\oplus L} $$
   $$ \boxtimes (i \hbar^{\nabla})|_{dx_m}^L, \boxplus (i \hbar^{\nabla}|_
   {dx_m}^L) $$



  $$ x^{{1 \over 2} + iy} = e^{x \log x} $$
  それゆえに、この定数はゼータ関数と微分幾何の量子化を因数にもつ
  素因子分解の式にもなっている。
  Therefore, this product also constructed with differential geometry
  of quantum level equation and zeta function.
  $$ = C $$ 
  そして、この関数はオイラーの定数から広中平祐定理による4重帰納法の
  オイラーの公式からの多様体積分へとのサーストン空間のスペクトラム関数
  ともなっている。
  And this function of Euler product respectrum of focus with
  Heisuke Hironaka manifold in four assembled of integral Euler equation.. 

  $$ C =  b_0 + \cfrac{c_1}{b_1 +
	  \cfrac{c_2}{b_2 +
	  \cfrac{c_3}{b_3 +
	  \cfrac{c_4}{b_4 + \cdots}}}} $$
  この方程式は指数による連分数としての役割も担っている。  
  This equation demanded with continued number of step function.
  $$ x^{{1 \over 2} + iy} = e^{x \log x} $$
  $$ (2.71828)^{2.828} = a^{a + b^{b + c^{c + d^{d \cdots }}}} $$
  $$ = \int e^f \cdot x^{1 - t}dx $$
  This represent is Gamma function in Euler product.
  Therefore this product is zeta function of global differential equation.

  $$ {d \over d f}F(v_{ij},h) = \int {e^{-f}\mathrm{dV}}
  [- \Delta v + \nabla_{i}\nabla_{j}
  v_{ij} - R_{ij}v_{ij} - v_{ij}\nabla_{i}\nabla_{j} + 2 <\nabla f,\nabla h>
  + (R + \nabla f^2)({v \over 2} - h)]$$
  $$ = {d \over d f}F = {2 {\int (R + \nabla_{i}\nabla_{j} f)^2} \over 
  { - (R + \Delta f)}}\mathrm{dm} $$
  これらの方程式は8種類の微分幾何の次元多様体として、そして、これらの
  多様体は曲平面による双対性をも生成している。 
  そして、このガウスの曲平面は、大域的微分多様体と微分幾何の量子化から
  素因数を形成してもいる。
  These equation are eight differential geometry of dimension calyement.
  And these calyement equation excluded into pair of dimension surface.
  This surface of Gauss function are global differential manifold,
  and differential geometry of quantum level.
  $$ F \ge {d \over d f}\int \int{1 \over (x \log x)^2}dx_m 
  + {d \over d f}\int \int{1 \over (y \log y)^{1 \over 2}}dy_m $$

  微分幾何の量子化はオイラーの定数とガンマ関数が指数による連分数
  としての不変性として素因数を形成していて、このガウスの曲平面
  による量子力学における重力場理論は、ダランベルシアンの切断多様体
  がこの大域的切断多様体を付加してもいる。
  Differential geometry of quantum level constructed with Euler product
  and Gamma function being discatastrophed from continued fraction style.
  Gravity equation lend with varint of monotonicity of level expresented
  from gravity of letter varient formula.
  これらの方程式は基本群と大域的微分多様体をエスコートしていもいて、
  ヴェイユ予想がこのダランベルシアンの切断方程式たちから
  輸送のポートにもなっている。ベータ関数とガンマ関数がこれらのフォームラ
  の方程式を放出してもいて、結果、これらの方程式は広中平祐定理の複素
  多様体とグリーシャ教授によるペレルマン多様体からサポートされてもいる。
  この2名の教授は、一つは抽象理論をもう一方は具象理論を説明としている。
  These equation escourted into Global differential manifold and
  fundemental equation.
  Weil's theorem is imported from this equation in gravity of letters.
  Beta function and Gamma function are excluded with these formulas.
  These equation comontius from Heisuke Hironaga of complex manifold
  and Gresha professor of Perelman manifold.
  These two professos are one of abstract theorem and the other of
  visual manifold theorem.


\end{document}
